\documentclass[12pt,a4paper]{article}


% PAQUETES

\usepackage[utf8]{inputenc}
\usepackage[T1]{fontenc}
\usepackage[spanish]{babel}
\usepackage{geometry}
\usepackage{array}
\usepackage{tabularx}
\usepackage{booktabs}
\usepackage{enumitem}
\usepackage{float}
\usepackage{longtable}
\usepackage{hyperref}
\usepackage{float}
\usepackage{placeins}

\geometry{margin=2.5cm}
\renewcommand{\arraystretch}{1.3}

\newcolumntype{Y}{>{\raggedright\arraybackslash}X}

\setlength{\parindent}{0pt}
\setlength{\parskip}{6pt}


% DOCUMENTO

\begin{document}

% PORTADA
\begin{titlepage}
    \centering

    {\large \textbf{UNIVERSIDAD TÉCNICA ESTATAL DE QUEVEDO}}\\[0.3cm]
    {\large Facultad de Ciencias de la Computación}\\
    {\large Carrera de Software}\\[1.5cm]

    {\Large \textbf{INGENIERÍA DE REQUERIMIENTOS}}\\[0.5cm]
    {\large ISR-401}\\[1.5cm]

    {\LARGE \textbf{FORM-PC-01}}\\[0.9cm]

    \rule{0.85\textwidth}{0.5pt}\\[0.5cm]

    {\large \textbf{Comparación de modelos de proceso de requisitos}}\\

    \rule{0.85\textwidth}{0.5pt}\\[1.5cm]

    \begin{flushleft}
        \textbf{Integrantes:}
        \begin{itemize}[noitemsep,topsep=2pt]
            \item Basurto Anchaguano Jordan Steven
            \item Fuertes Calvo Landy Lorena
            \item Llerena Abril Angelina Julexy
            \item Monserrate Ochoa Erick Ariel
        \end{itemize}
        \vspace{0.2cm}
        \textbf{Nivel:} Cuarto\\[0.4cm]
        \textbf{Paralelo:} A\\[0.4cm]
        \textbf{Docente:} Gleiston Cicerón Guerrero Ulloa, PhD\\[0.4cm]
        \textbf{Modalidad:} Grupal\\[0.4cm]
        \textbf{Fecha de entrega:} 7 de noviembre de 2026
    \end{flushleft}

    \vfill

    {\large Quevedo -- Ecuador}\\
    {\large 2026}

\end{titlepage}

\section{Introducción}

El presente trabajo analiza la aplicación de los modelos
secuencial, iterativo e incremental y ágil en cuatro
escenarios de proyectos de software. Para ello, se
consideran las características de cada organización,
la estabilidad de los requisitos, la disponibilidad de
las partes interesadas, los riesgos y las condiciones
normativas o contractuales.

Además, se selecciona un modelo de proceso para cada
escenario y se aplica el framework de Ingeniería de
Requisitos de Pohl a uno de ellos, identificando sus
actores, actividades centrales, actividades
transversales y artefactos.

\section{Criterios de comparación}

Para comparar los tres modelos de proceso se utilizan
los siguientes criterios:

\begin{enumerate}[label=\textbf{K\arabic*.}]
    \item \textbf{Estabilidad y gestión de cambios:}
    analiza cuánto y cuándo se espera que cambien los
    requisitos del sistema.

    \item \textbf{Partes interesadas y disponibilidad:}
    identifica quiénes participan en el proyecto, su
    disponibilidad y quién tiene autoridad para tomar
    decisiones.

    \item \textbf{Tipo de sistema y criticidad:}
    considera el tipo de sistema y las consecuencias
    que podrían generar los errores en su funcionamiento.

    \item \textbf{Riesgo principal del proyecto:}
    identifica las principales amenazas técnicas,
    organizativas o de plazo que podrían afectar el
    desarrollo.

    \item \textbf{Documentación y marco normativo o
    contractual:}
    considera los documentos, normas y condiciones
    externas que debe cumplir el proyecto.
\end{enumerate}


% ESCENARIO 1

\section{Escenario 1. Control de temperatura y nivel en una extractora de aceite de palma}

\subsection{Descripción del escenario}

\subsection{Matriz comparativa}

\begin{table}[H]
\centering
\caption{Matriz comparativa del escenario 1}
\small
\begin{tabularx}{\textwidth}{
    >{\raggedright\arraybackslash}p{2.4cm} Y Y Y
}
\toprule
\textbf{Criterio} &
\textbf{Secuencial (cascada)} &
\textbf{Iterativo e incremental} &
\textbf{Ágil} \\
\midrule

\textbf{K1. Estabilidad y cambios} & & & \\
\midrule
\textbf{K2. Partes interesadas y disponibilidad} & & & \\
\midrule
\textbf{K3. Tipo de sistema y criticidad} & & & \\
\midrule
\textbf{K4. Riesgo principal} & & & \\
\midrule
\textbf{K5. Documentación y marco normativo o contractual} & & & \\
\bottomrule
\end{tabularx}
\end{table}

\subsection{Recomendación y riesgos}

\begin{table}[H]
\centering
\caption{Recomendación y riesgos del escenario 1}
\small
\begin{tabularx}{\textwidth}{
    >{\raggedright\arraybackslash}p{3.5cm} Y
}
\toprule
\textbf{Elemento} & \textbf{Respuesta} \\
\midrule
Modelo recomendado & \\
\midrule
Razón 1 & \\
\midrule
Razón 2 & \\
\midrule
Razón 3 & \\
\midrule
Por qué se descartan los otros dos modelos & \\
\midrule
Tensión principal del escenario & \\
\midrule
Riesgo 1 y mitigación & \\
\midrule
Riesgo 2 y mitigación & \\
\bottomrule
\end{tabularx}
\end{table}

\subsection{Actores y framework de Ingeniería de Requisitos}

\subsubsection{Identificación de actores}

\begin{table}[H]
\centering
\caption{Actores del escenario 1}
\small
\begin{tabularx}{\textwidth}{
    >{\raggedright\arraybackslash}p{2.8cm} Y Y
}
\toprule
\textbf{Actor del escenario} &
\textbf{Rol en el proceso de requisitos} &
\textbf{Cuándo y cómo participa} \\
\midrule
& & \\
\midrule
& & \\
\midrule
& & \\
\midrule
& & \\
\midrule
& & \\
\bottomrule
\end{tabularx}
\end{table}

\subsubsection{Actividades centrales, transversales y artefactos}

\begin{table}[H]
\centering
\caption{Aplicación del framework de IR en el escenario 1}
\scriptsize
\begin{tabularx}{\textwidth}{
    >{\raggedright\arraybackslash}p{2.1cm} Y Y Y
}
\toprule
\textbf{Fase o iteración} &
\textbf{Actividades centrales} &
\textbf{Actividades transversales} &
\textbf{Artefactos producidos} \\
\midrule
& & & \\
\midrule
& & & \\
\midrule
& & & \\
\midrule
& & & \\
\bottomrule
\end{tabularx}
\end{table}


% ESCENARIO 2

\section{Escenario 2. Aplicación móvil de pedidos para un emprendimiento de comida típica}

\subsection{Descripción del escenario}

Sabores de la Abuela Rosa es un emprendimiento familiar de Quevedo que vende
bolones, tigrillo, encebollado y tortillas de yuca en un local pequeño del
centro de la ciudad. Recibe unos 70 pedidos diarios por una aplicación de
mensajería; la dueña, Doña Rosa Mendoza, los anota en un cuaderno y con
frecuencia se equivoca de dirección o pierde mensajes en las horas pico.

Se requiere una aplicación móvil para que los clientes vean el menú del día,
hagan pedidos, elijan retirar en el local o recibir a domicilio y paguen por
transferencia (más adelante, quizá, con tarjeta). Doña Rosa responde a
cualquier hora y prueba cada versión, pero cambia de idea con facilidad; su
hijo propone funciones nuevas en cada reunión y el menú cambia casi a diario.
La aplicación guardará nombres, direcciones y teléfonos de clientes, por lo
que debe cumplir la Ley Orgánica de Protección de Datos Personales (LOPDP), y
los pagos por transferencia deben conciliarse para evitar pedidos sin pagar.

El presupuesto es fijo y pequeño, el plazo es de 14 semanas con fecha fija
(fiestas de la ciudad) y el equipo está formado por dos desarrolladores a
medio tiempo. Muchos clientes usan teléfonos Android de gama baja con datos
móviles intermitentes, por lo que la aplicación debe ser liviana y tolerar
cortes de conexión.

\subsection{Matriz comparativa}

\begin{table}[H]
\centering
\caption{Matriz comparativa del escenario 2}
\small
\begin{tabularx}{\textwidth}{
    >{\raggedright\arraybackslash}p{2.4cm} Y Y Y
}
\toprule
\textbf{Criterio} &
\textbf{Secuencial (cascada)} &
\textbf{Iterativo e incremental} &
\textbf{Ágil} \\
\midrule

\textbf{K1. Estabilidad y cambios} & \textbf{Baja adecuación.} Exige congelar los requisitos al inicio; con un menú diario y una dueña que cambia de idea, cada cambio obligaría a rehacer fases. & \textbf{Media.} Absorbe cambios entre iteraciones, pero sobre un plan de iteraciones definido de antemano que aquí quedaría obsoleto. & \textbf{Alta.} El backlog se reprioriza en cada sprint; los cambios se esperan y se gestionan sin romper el trabajo en curso. \\
\midrule
\textbf{K2. Partes interesadas y disponibilidad} & \textbf{Baja.} Desaprovecha a una dueña disponible a diario: solo participa al inicio y en la entrega final. & \textbf{Media-alta.} La dueña valida al final de cada iteración, aunque con menor frecuencia de la que puede ofrecer. & \textbf{Alta.} Doña Rosa actúa como Product Owner: decide, prioriza las ideas de su hijo y prueba cada incremento. \\
\midrule
\textbf{K3. Tipo de sistema y criticidad} & \textbf{Media.} Su control formal es excesivo para una app comercial pequeña sin riesgo para la vida. & \textbf{Alta.} Permite revisar pagos y datos personales en cada iteración. & \textbf{Alta.} Suficiente si se fijan criterios de aceptación y pruebas para pagos por transferencia y datos de clientes. \\
\midrule
\textbf{K4. Riesgo principal} & \textbf{Baja.} Entrega todo al final: si hay retraso no habrá nada para las fiestas, y el rendimiento en Android de gama baja se descubre tarde. & \textbf{Media-alta.} Hay versiones parciales utilizables y pruebas en dispositivos reales por iteración. & \textbf{Alta.} Fecha y costo fijos con alcance variable: siempre existe un MVP usable; la conectividad se prueba desde el primer mes. \\
\midrule
\textbf{K5. Documentación y marco normativo o contractual} & \textbf{Media-alta.} Su documentación facilita evidenciar el cumplimiento de la LOPDP, pero pesa demasiado para dos personas a medio tiempo. & \textbf{Media-alta.} Permite documentar por iteración con un esfuerzo moderado. & \textbf{Media.} Exige incluir explícitamente la LOPDP y la conciliación de pagos en la Definición de Hecho; el acuerdo informal de precio fijo se respalda con una priorización MoSCoW escrita. \\
\bottomrule
\end{tabularx}
\end{table}

\subsection{Recomendación y riesgos}

\begin{table}[H]
\centering
\caption{Recomendación y riesgos del escenario 2}
\small
\begin{tabularx}{\textwidth}{
    >{\raggedright\arraybackslash}p{3.5cm} Y
}
\toprule
\textbf{Elemento} & \textbf{Respuesta} \\
\midrule
Modelo recomendado & Ágil (Scrum ligero): 7 sprints de 2 semanas, con fecha y presupuesto fijos y alcance variable priorizado como producto mínimo viable (MVP). \\
\midrule
Razón 1 & La volatilidad de los requisitos es muy alta (menú diario, ideas nuevas de la dueña y su hijo); el backlog priorizado absorbe esos cambios sin rehacer fases. \\
\midrule
Razón 2 & La dueña está disponible a diario y prueba cada versión, condición ideal para el rol de Product Owner y la validación continua. \\
\midrule
Razón 3 & Con fecha fija y presupuesto cerrado, entregar por incrementos garantiza un MVP usable en la semana 8 y una versión estable en la semana 14; además, el equipo pequeño no puede sostener documentación pesada. \\
\midrule
Por qué se descartan los otros dos modelos & Cascada: requiere requisitos estables, entrega todo al final y su documentación consume tiempo que el equipo no tiene. Iterativo e incremental: es viable, pero planifica las iteraciones de antemano y el contenido cambia demasiado para ese plan. \\
\midrule
Tensión principal del escenario & Doña Rosa quiere ``tener todo listo'' para una fecha fija con presupuesto cerrado, pero cambia de idea con facilidad: alcance fijo y volatilidad alta no son compatibles, por lo que se fija la fecha y se negocia el alcance. \\
\midrule
Riesgo 1 y mitigación & \textbf{Riesgo:} el alcance crece con cada idea nueva (puntos, promociones, propuestas del hijo) y compromete la fecha. \textbf{Mitigación:} toda idea entra al backlog y nunca al sprint en curso; priorización MoSCoW acordada por escrito con la dueña. \\
\midrule
Riesgo 2 y mitigación & \textbf{Riesgo:} incumplir la LOPDP o aceptar pedidos sin pago. \textbf{Mitigación:} aviso de privacidad, consentimiento explícito, cifrado y eliminación de cuenta como historias obligatorias desde el sprint 1; ningún pedido pasa a cocina sin comprobante conciliado; ambos incluidos en la Definición de Hecho. \\
\bottomrule
\end{tabularx}
\end{table}

\subsection{Actores y framework de Ingeniería de Requisitos}

\subsubsection{Identificación de actores}

\begin{table}[H]
\centering
\caption{Actores del escenario 2}
\small
\begin{tabularx}{\textwidth}{
    >{\raggedright\arraybackslash}p{2.8cm} Y Y
}
\toprule
\textbf{Actor del escenario} &
\textbf{Rol en el proceso de requisitos} &
\textbf{Cuándo y cómo participa} \\
\midrule
Doña Rosa Mendoza (dueña) & Product Owner y principal fuente de requisitos del negocio; decide el alcance y prioriza el backlog. & Durante todo el proyecto: entrevistas y talleres al inicio, revisión y prueba de cada incremento al final de cada sprint, respuestas diarias por mensajería. \\
\midrule
Hijo de Doña Rosa (redes sociales) & Stakeholder que propone funciones de promoción y fidelización. & En las reuniones de revisión; sus ideas se registran en el backlog y se priorizan con la dueña, nunca dentro del sprint en curso. \\
\midrule
Clientes del emprendimiento & Usuarios finales; fuente de requisitos de usabilidad, rendimiento y conectividad. & Encuesta corta al inicio, pruebas con 3 a 5 clientes frecuentes en el sprint de pedidos y piloto real del MVP en la semana 8. \\
\midrule
Desarrolladores (2, medio tiempo) & Equipo que elicita, analiza, documenta y valida los requisitos; uno actúa además como Scrum Master. & En todos los sprints: planificación, estimación, construcción, demostraciones y mantenimiento de la trazabilidad. \\
\midrule
Personal de cocina, despacho y repartidor & Usuarios internos del panel de pedidos y de la información de entrega. & Observación de una hora pico al inicio, entrevista sobre direcciones y validación del flujo de pedidos y pagos antes del piloto. \\
\bottomrule
\end{tabularx}
\end{table}

\subsubsection{Actividades centrales, transversales y artefactos}

\begin{table}[H]
\centering
\caption{Aplicación del framework de IR en el escenario 2}
\scriptsize
\begin{tabularx}{\textwidth}{
    >{\raggedright\arraybackslash}p{2.1cm} Y Y Y
}
\toprule
\textbf{Fase o iteración} &
\textbf{Actividades centrales} &
\textbf{Actividades transversales} &
\textbf{Artefactos producidos} \\
\midrule
Sprint 0 (semanas 1--2): visión y backlog inicial & \textbf{Elicitación:} entrevista con Doña Rosa y su hijo, observación de una hora pico, revisión del cuaderno y los chats, encuesta a clientes. \textbf{Negociación:} separar necesidades de deseos y acordar el MVP. \textbf{Documentación:} visión, historias de usuario y requisitos no funcionales (LOPDP, peso de la app, modo sin conexión). & \textbf{Validación:} la dueña aprueba el MVP y la priorización MoSCoW por escrito. \textbf{Gestión:} backlog como línea base, matriz de trazabilidad y acuerdo de comunicación (revisión quincenal, canal único de ideas). & Visión del producto, mapa de historias, backlog priorizado con criterios de aceptación, matriz de trazabilidad, acta MoSCoW. \\
\midrule
Sprints 1--3 (semanas 3--8): menú, pedidos y pago por transferencia (MVP) & \textbf{Elicitación:} talleres sobre el menú diario y la verificación de transferencias; entrevista al repartidor; pruebas con 3 a 5 clientes. \textbf{Negociación:} edición simple del menú desde el teléfono, referencia y GPS opcional en direcciones, estado ``pendiente de pago''. \textbf{Documentación:} historias de menú, carrito, entrega, comprobante y conciliación. & \textbf{Validación:} demos en Android de gama baja, pruebas de aceptación con pedidos simulados en hora pico y piloto real del MVP durante una semana. \textbf{Gestión:} ideas nuevas registradas y pospuestas; trazabilidad de los errores del cuaderno a las historias; revisión LOPDP del MVP. & Prototipos en papel o Figma, historias con criterios de aceptación, Definición de Hecho, incremento MVP, informe del piloto. \\
\midrule
Sprints 4--5 (semanas 9--12): tolerancia a cortes, rendimiento y fidelización & \textbf{Elicitación:} métricas del piloto (fallos de conexión, tiempos de carga, quejas); taller sobre puntos o promociones. \textbf{Negociación:} priorizar el modo sin conexión y elegir una sola opción de fidelización según presupuesto y plazo. \textbf{Documentación:} requisitos no funcionales medibles (p. ej., app menor a 15 MB, menú visible sin conexión) e historias de la opción elegida. & \textbf{Validación:} pruebas en teléfonos de gama baja con datos intermitentes; revisión del sprint con la familia y algunos clientes. \textbf{Gestión:} repriorización del backlog con datos del piloto; la opción descartada queda registrada para después de las fiestas. & Backlog actualizado, requisitos no funcionales medibles, matriz de trazabilidad actualizada, incremento con modo sin conexión y notificaciones. \\
\midrule
Sprint 6 (semanas 13--14): estabilización para las fiestas & \textbf{Elicitación:} revisión de defectos reportados y previsión del doble de pedidos. \textbf{Negociación:} congelar el alcance; no entran funciones nuevas. \textbf{Documentación:} ajustes finales, manual breve para la dueña y aviso de privacidad definitivo. & \textbf{Validación:} prueba de carga con el doble de pedidos y aceptación final de Doña Rosa. \textbf{Gestión:} línea base final de requisitos y backlog posterior a las fiestas (pago con tarjeta, opción de fidelización descartada). & Versión estable publicada, manual de uso, aviso de privacidad, acta de aceptación, backlog para la siguiente etapa. \\
\bottomrule
\end{tabularx}
\end{table}


% ESCENARIO 3

\section{Escenario 3. Sistema de trámites municipales contratado por contratación pública}

\subsection{Descripción del escenario}

\subsection{Matriz comparativa}

\begin{table}[H]
\centering
\caption{Matriz comparativa del escenario 3}
\small
\begin{tabularx}{\textwidth}{
    >{\raggedright\arraybackslash}p{2.4cm} Y Y Y
}
\toprule
\textbf{Criterio} &
\textbf{Secuencial (cascada)} &
\textbf{Iterativo e incremental} &
\textbf{Ágil} \\
\midrule

\textbf{K1. Estabilidad y cambios} & & & \\
\midrule
\textbf{K2. Partes interesadas y disponibilidad} & & & \\
\midrule
\textbf{K3. Tipo de sistema y criticidad} & & & \\
\midrule
\textbf{K4. Riesgo principal} & & & \\
\midrule
\textbf{K5. Documentación y marco normativo o contractual} & & & \\
\bottomrule
\end{tabularx}
\end{table}

\subsection{Recomendación y riesgos}

\begin{table}[H]
\centering
\caption{Recomendación y riesgos del escenario 3}
\small
\begin{tabularx}{\textwidth}{
    >{\raggedright\arraybackslash}p{3.5cm} Y
}
\toprule
\textbf{Elemento} & \textbf{Respuesta} \\
\midrule
Modelo recomendado & \\
\midrule
Razón 1 & \\
\midrule
Razón 2 & \\
\midrule
Razón 3 & \\
\midrule
Por qué se descartan los otros dos modelos & \\
\midrule
Tensión principal del escenario & \\
\midrule
Riesgo 1 y mitigación & \\
\midrule
Riesgo 2 y mitigación & \\
\bottomrule
\end{tabularx}
\end{table}

\subsection{Actores y framework de Ingeniería de Requisitos}

\subsubsection{Identificación de actores}

\begin{table}[H]
\centering
\caption{Actores del escenario 3}
\small
\begin{tabularx}{\textwidth}{
    >{\raggedright\arraybackslash}p{2.8cm} Y Y
}
\toprule
\textbf{Actor del escenario} &
\textbf{Rol en el proceso de requisitos} &
\textbf{Cuándo y cómo participa} \\
\midrule
& & \\
\midrule
& & \\
\midrule
& & \\
\midrule
& & \\
\midrule
& & \\
\bottomrule
\end{tabularx}
\end{table}

\subsubsection{Actividades centrales, transversales y artefactos}

\begin{table}[H]
\centering
\caption{Aplicación del framework de IR en el escenario 3}
\scriptsize
\begin{tabularx}{\textwidth}{
    >{\raggedright\arraybackslash}p{2.1cm} Y Y Y
}
\toprule
\textbf{Fase o iteración} &
\textbf{Actividades centrales} &
\textbf{Actividades transversales} &
\textbf{Artefactos producidos} \\
\midrule
& & & \\
\midrule
& & & \\
\midrule
& & & \\
\midrule
& & & \\
\bottomrule
\end{tabularx}
\end{table}

% ESCENARIO 4

\section{Escenario 4. Plataforma web para una cooperativa de productores de cacao con socios dispersos}

\subsection{Descripción del escenario}

La Cooperativa Agroproductiva Frutos del Estero (organización ficticia creada para esta actividad) agrupa
a 640 pequeños productores de cacao de recintos rurales de Mocache, Ventanas y Valencia. La cooperativa acopia el cacao en un centro ubicado en Mocache, lo vende a exportadores y paga a cada
socio según el peso, la humedad y la calidad de su entrega. Hoy las entregas se anotan en tarjetas de
papel, las liquidaciones se calculan en una hoja de cálculo y los socios se quejan de que no entienden
cómo se calculó su pago ni cuánto deben de los pequeños créditos para insumos. La directiva quiere
una plataforma web donde los socios consulten entregas, liquidaciones y saldos de crédito, y donde la
cooperativa publique avisos y convocatorias a asambleas.

La gerente, Ing. Patricia Vera (persona ficticia), está disponible casi a diario, pero no puede decidir sola:
los cambios importantes los aprueba la directiva en su reunión mensual. Los socios viven lejos y solo se
les puede consultar los días de acopio o en las asambleas trimestrales; muchos tienen poca experiencia
con tecnología. Las reglas de liquidación están en el estatuto y son estables, pero los compradores exportadores cambian cada temporada los requisitos de trazabilidad y certificación del cacao, y la plataforma
deberá registrarlos. Se manejará dinero de los socios y datos personales, y la cooperativa está sujeta
al control del organismo que supervisa a la economía popular y solidaria. El proyecto se financia con
una donación que exige informes semestrales de avance con funciones terminadas. El equipo es de 4
personas y el plazo es de 12 meses. En el centro de acopio la conexión a internet es inestable, por lo que
el registro de entregas debe funcionar sin conexión y sincronizarse después.

\FloatBarrier
\subsection{Matriz comparativa}

\begin{table}[H]
\centering
\caption{Matriz comparativa del escenario 4}
\label{tab:matriz_escenario4}
\footnotesize
\renewcommand{\arraystretch}{1.05}
\setlength{\tabcolsep}{3pt}
\begin{tabularx}{\textwidth}{
    >{\raggedright\arraybackslash}p{2.35cm} Y Y Y
}
\toprule
\textbf{Criterio} &
\textbf{Secuencial (cascada)} &
\textbf{Iterativo e incremental} &
\textbf{Ágil} \\
\midrule

\textbf{K1. Estabilidad y cambios}
&
Poco conveniente, porque los requisitos de trazabilidad y certificación cambian cada temporada y podrían modificar lo definido al inicio.
&
Conveniente, porque las reglas de liquidación pueden establecerse inicialmente y los requisitos de trazabilidad actualizarse en ciclos posteriores.
&
Conveniente, porque permite incorporar cambios frecuentes de los compradores mediante la actualización continua de prioridades.
\\
\midrule

\textbf{K2. Partes interesadas y disponibilidad}
&
Conveniente en parte, porque la gerente está disponible casi a diario y la directiva puede aprobar requisitos en reuniones mensuales. Sin embargo, los socios participan pocas veces.
&
Conveniente, porque permite revisar cada incremento con la gerente y aprovechar los días de acopio y las asambleas para validar las necesidades de los socios.
&
Poco conveniente, porque la participación frecuente de los socios es difícil y la gerente no tiene autoridad para aprobar por sí sola los cambios importantes.
\\
\midrule

\textbf{K3. Tipo de sistema y criticidad}
&
Conveniente en parte, porque el manejo de dinero, créditos y datos personales requiere requisitos definidos y revisiones cuidadosas antes de la implementación.
&
Muy conveniente, porque permite implementar y comprobar progresivamente las funciones de entregas, liquidaciones y créditos antes de ampliar el sistema.
&
Conveniente en parte, porque facilita adaptar las funcionalidades, pero los cambios frecuentes necesitan controles para evitar errores en los pagos y datos de los socios.
\\
\midrule

\textbf{K4. Riesgo principal}
&
Poco conveniente, porque los problemas del funcionamiento sin conexión y los cambios de trazabilidad podrían descubrirse tarde, cuando modificar el sistema resulte más costoso.
&
Muy conveniente, porque permite probar desde los primeros ciclos el registro sin conexión, la sincronización y las funcionalidades principales.
&
Conveniente, porque permite recibir retroalimentación y ajustar funcionalidades, aunque la poca disponibilidad de los socios dificulta validar continuamente los cambios.
\\
\midrule

\textbf{K5. Documentación y marco normativo o contractual}
&
Conveniente, porque facilita establecer documentación formal de los requisitos y mantener registros de las decisiones relacionadas con el manejo del dinero y los datos personales.
&
Muy conveniente, porque permite mantener documentación actualizada por versiones y presentar funciones terminadas en los informes semestrales del donante.
&
Conveniente en parte, porque permite documentar los cambios por entregas, pero debe evitarse que la documentación sea insuficiente para las necesidades de control del proyecto.
\\
\bottomrule
\end{tabularx}
\end{table}

\FloatBarrier
\subsection{Recomendación y riesgos}

\begin{table}[H]
\centering
\caption{Recomendación y riesgos del escenario 4}
\label{tab:recomendacion_escenario4}
\small
\renewcommand{\arraystretch}{1.1}
\setlength{\tabcolsep}{4pt}
\begin{tabularx}{\textwidth}{
    >{\raggedright\arraybackslash}p{3.5cm} Y
}
\toprule
\textbf{Elemento} & \textbf{Respuesta} \\
\midrule

Modelo recomendado
&
Iterativo e incremental.
\\
\midrule

Razón 1 (K1)
&
Las reglas de liquidación están establecidas en el estatuto y son estables, mientras que los requisitos de trazabilidad y certificación cambian cada temporada. Por ello, se pueden desarrollar inicialmente las funciones de liquidación y actualizar las de trazabilidad en ciclos posteriores.
\\
\midrule

Razón 2 (K2)
&
La gerente está disponible casi a diario, pero los cambios importantes necesitan aprobación de la directiva en sus reuniones mensuales. El modelo iterativo e incremental permite organizar las revisiones según la disponibilidad de cada participante, sin depender de reuniones continuas.
\\
\midrule

Razón 3 (K5)
&
El donante solicita informes semestrales con funciones terminadas. Los incrementos permiten presentar avances verificables y mantener la documentación de los requisitos actualizada durante el desarrollo.
\\
\midrule

Por qué se descartan los otros dos modelos
&
\textbf{Secuencial (cascada):} se descarta porque los requisitos de trazabilidad y certificación cambian cada temporada, lo que podría generar modificaciones costosas si todo se define al inicio.

\medskip

\textbf{Ágil:} se descarta porque los socios tienen poca disponibilidad para participar continuamente y la gerente no puede aprobar sola los cambios importantes.
\\
\midrule

Tensión principal del escenario
&
Aunque el modelo iterativo e incremental es el más conveniente, los cambios frecuentes en los requisitos de trazabilidad y certificación de los compradores apuntan hacia el modelo ágil, debido a que requieren adaptaciones continuas.
\\
\midrule

Riesgo 1 y mitigación
&
\textbf{Riesgo:} los cambios en los requisitos de trazabilidad podrían no identificarse a tiempo entre iteraciones.

\medskip

\textbf{Mitigación:} consultar periódicamente los nuevos requisitos de los compradores y actualizar los requisitos antes de iniciar cada iteración.
\\
\midrule

Riesgo 2 y mitigación
&
\textbf{Riesgo:} la sincronización de datos podría generar inconsistencias al trabajar sin conexión.

\medskip

\textbf{Mitigación:} realizar pruebas de funcionamiento sin conexión y sincronización en cada incremento antes de implementarlo.
\\
\bottomrule
\end{tabularx}
\end{table}

\FloatBarrier
\subsection{Actores y framework de Ingeniería de Requisitos}

\subsubsection{Identificación de actores}

Se identifican las partes interesadas que intervienen directamente en el proceso de requisitos y aquellas externas que establecen condiciones que deben considerarse durante el desarrollo.

\begin{table}[H]
\centering
\caption{Actores del escenario 4}
\label{tab:actores_escenario4}
\small
\renewcommand{\arraystretch}{1.1}
\setlength{\tabcolsep}{4pt}
\begin{tabularx}{\textwidth}{
    >{\raggedright\arraybackslash}p{2.8cm} Y Y
}
\toprule
\textbf{Actor del escenario} &
\textbf{Rol en el proceso de requisitos} &
\textbf{Cuándo y cómo participa} \\
\midrule

Gerente (Ing. Patricia Vera)
&
\textbf{Directo.} Representa a la cooperativa y comunica sus necesidades y prioridades.
&
Participa casi a diario, proporciona información y revisa los avances del proyecto.
\\
\midrule

Directiva
&
\textbf{Directo.} Aprueba los cambios importantes y toma decisiones sobre el proyecto.
&
Participa en reuniones mensuales para revisar y aprobar los cambios que correspondan.
\\
\midrule

Socios productores
&
\textbf{Directos.} Son los usuarios finales de la plataforma y expresan sus necesidades de consulta.
&
Participan durante los días de acopio y las asambleas trimestrales mediante consultas y pruebas de las funcionalidades.
\\
\midrule

Equipo del proyecto
&
\textbf{Directo.} Analiza, documenta, desarrolla y prueba los requisitos.
&
Participa durante todo el desarrollo mediante actividades de elicitación, documentación, validación y gestión de requisitos.
\\
\midrule

Compradores exportadores
&
\textbf{Indirectos.} Establecen los requisitos de trazabilidad y certificación del cacao.
&
Sus requisitos se consultan en cada temporada para actualizar las necesidades del sistema.
\\
\midrule

Organismo supervisor
&
\textbf{Indirecto.} Establece condiciones de control que la cooperativa debe considerar.
&
Sus disposiciones se toman en cuenta al definir y revisar los requisitos de la plataforma.
\\
\midrule

Donante
&
\textbf{Indirecto.} Establece condiciones para el seguimiento y financiamiento del proyecto.
&
Recibe informes semestrales de avance con funciones terminadas.
\\
\bottomrule
\end{tabularx}
\end{table}

\FloatBarrier
\subsubsection{Actividades centrales, transversales y artefactos}

Para aplicar el framework de Ingeniería de Requisitos de Pohl, se organizan las actividades centrales de elicitación, documentación, negociación y validación. De manera transversal, se considera la gestión de requisitos durante las cuatro iteraciones propuestas.

\begin{table}[H]
\centering
\caption{Aplicación del framework de IR en el escenario 4}
\label{tab:framework_escenario4}
\scriptsize
\renewcommand{\arraystretch}{1.05}
\setlength{\tabcolsep}{3pt}
\begin{tabularx}{\textwidth}{
    >{\raggedright\arraybackslash}p{2.1cm} Y Y Y
}
\toprule
\textbf{Fase o iteración} &
\textbf{Actividades centrales} &
\textbf{Actividades transversales} &
\textbf{Artefactos producidos} \\
\midrule

\textbf{Iteración 1: Análisis inicial}
&
\textbf{Elicitación:} recopilar las necesidades de la gerente, la directiva y los socios mediante entrevistas y consultas durante los días de acopio y las asambleas.

\textbf{Documentación:} registrar las funciones principales de la plataforma y los requisitos iniciales.

\textbf{Negociación:} acordar las prioridades y el alcance inicial con la gerente y la directiva.
&
\textbf{Validación:} revisar los requisitos con la gerente y gestionar su aprobación con la directiva.

\textbf{Gestión:} establecer una primera versión de los requisitos y registrar las decisiones y cambios.
&
Metas del sistema, escenarios de consulta, requisitos iniciales y registro de decisiones.
\\
\midrule

\textbf{Iteración 2: Entregas y liquidaciones}
&
\textbf{Elicitación:} recopilar las necesidades para registrar entregas, calcular liquidaciones y consultar créditos.

\textbf{Documentación:} especificar los requisitos relacionados con el peso, la humedad, la calidad y los cálculos de pago.

\textbf{Negociación:} acordar los criterios de registro y consulta con la cooperativa.
&
\textbf{Validación:} revisar los cálculos de liquidación con la gerente y algunos socios productores para comprobar que los resultados sean comprensibles.

\textbf{Gestión:} registrar los cambios y mantener actualizados los requisitos.
&
Escenarios de registro de entregas, liquidaciones y consulta de créditos; requisitos funcionales y criterios de aceptación.
\\
\midrule

\textbf{Iteración 3: Trazabilidad y funcionamiento sin conexión}
&
\textbf{Elicitación:} identificar los requisitos de trazabilidad, certificación, registro sin conexión y sincronización.

\textbf{Documentación:} especificar las condiciones de funcionamiento sin conexión y actualización posterior de los datos.

\textbf{Negociación:} acordar los cambios necesarios según los requisitos de los compradores.
&
\textbf{Validación:} comprobar los requisitos de trazabilidad con las condiciones comunicadas por los compradores y probar el registro y la sincronización.

\textbf{Gestión:} controlar las modificaciones y mantener las versiones actualizadas.
&
Escenarios de trazabilidad, requisitos de funcionamiento sin conexión, requisitos de sincronización y especificaciones actualizadas.
\\
\midrule

\textbf{Iteración 4: Consultas y cierre}
&
\textbf{Elicitación:} recopilar las necesidades relacionadas con consultas, avisos y convocatorias de asambleas.

\textbf{Documentación:} completar los requisitos de las funcionalidades y las condiciones de uso de la plataforma.

\textbf{Negociación:} revisar los requisitos pendientes y acordar cuáles se incorporarán al alcance final.
&
\textbf{Validación:} realizar pruebas de consulta y navegación con socios productores, considerando su poca experiencia tecnológica, y revisar las funcionalidades con la gerente.

\textbf{Gestión:} revisar los requisitos pendientes y mantener la especificación final actualizada.
&
Requisitos finales, escenarios de uso, resultados de validación y especificación actualizada del sistema.
\\
\bottomrule
\end{tabularx}
\end{table}

\FloatBarrier


% FIN DEL DOCUMENTO

\end{document}
